\section{Exact first-violation
macros}\label{exact-first-violation-macros}

The preceding decision can be strengthened to return the first event
itself and its complete post-event state. This requires a finite-path
argument inside the nonsimple component, rather than an assumption of
periodicity there.

\textbf{Theorem F (exact event and finite prefixes).} A complete finite
field may be given as the XOR of \(M\) finite arithmetic segments whose
strides belong to
\[\{(0,0),(\pm2,\pm2),(\pm2,0),(0,\pm2),(\pm4,0),(0,\pm4)\}.\] Zero
stride denotes a singleton. The first future \(J_{10}\), or its
permanent absence, is decidable in polynomial bit time in this
representation. A finite answer returns its exact time and full state
just after the event. One call uses at most five periodic plans and 1111
literal main-orbit updates, adding at most 1530 segment terms.

Let \(W_0\) be the diameter in the maximum norm of the actual initial
support and head. Conditional on \(n\) events existing, their post-event
time satisfies \[s_n\le n(212W_0+854199)+415308n(n-1). \tag{MACRO1}\]
Every intermediate complete state lies in the initial convex hull
enlarged by \([-1959n,1959n]^2\). The first \(n\) events and their
complete encoded endpoint are computable in
\(\operatorname{poly}(M+n+b)\) bit operations, with \(n\) measured
numerically and \(b\) the original coordinate and length bit size.

These macros preserve finite histories exactly. Their lift counts fold
starts within one event-free traversal; a fresh event begins a fresh
traversal. The following proof supplies both the finite graph coverage
and the arithmetic continuation needed for this stronger conclusion.

\subsection{Pinned base and finite-path coverage}\label{macro-3}

The base contains every physically feasible 150-letter word all of whose
ten-factors are in L10. Its 6605 edge words generate a graph G with 6501
vertices, the 149-letter prefixes/suffixes. Physical feasibility is
exactly alternation of requested turns on consecutive reads of the same
physical cell. This uses the established forbidden-word method
here.

Every actual event-free continuation, once its 149-turn history is a G
vertex, follows G. At most 149 nonviolating ordinary updates are needed
to obtain such a vertex: afterwards the last 149 turns are wholly new,
and if the next turn is legal their feasible 150-word is in the complete
enumeration. If the next turn is illegal, it is returned as the event
instead. This covers input time zero without inventing a periodic
negative-time past.

The pinned graph has simple W and T components of lengths 104 and 24,
and one component H with 366 vertices and 370 internal edges. The prior
explicit quotient maps H to phases 0,\ldots,101. There is an edge
(j,B\_j)-\textgreater j+1 mod 102, plus (0,L)-\textgreater75, where B is
a 102-letter fold word beginning R. The short loop is A=L
B\_75\ldots B\_101, length 28. Their literal geometry is

\begin{verbatim}
A: displacement (0,-4), heading restored;
B: displacement (1,11), heading reversed.
\end{verbatim}

Every internal H edge is checked against this quotient, not merely
sampled. The already supplied all-parameter lemma is

\begin{verbatim}
B A^n B A^m B is physically infeasible for ALL n,m>=0. (MACRO2)
\end{verbatim}

Its explicit consecutive same-cell/same-R conflicts and their parameter
coverage are in Subsection \ref{ten-4}. A finite word need not extend to
an infinite legal word for (MACRO2) to apply. This is the essential
finite-path rather than infinite-cycle input.

The base edge set and quotient have frozen SHA256 values in the source
witness. The event runtime regenerates both and checks exact equality.

\subsection{A complete bounded-fold lift}\label{macro-4}

Outside H keep each vertex v with counter 0. Inside H use (v,c),
c=0,1,2,3, except that (v,3) with quotient phase 0 is forbidden. On an
internal H edge increase c exactly when its quotient starts the long
loop, namely phase 0 with appended R. All other internal edges keep
c.~On entry into H start c=0; on exit forget c.~Disallow c\textgreater3
and disallow arrival at phase 0 with c=3.

This permits two complete B loops and a third incomplete B loop; a third
*completed* loop is excluded. An initial partial loop is not charged, so
no assumption is made about the path before H entry. Between three
long-loop starts at phase zero there are powers of A, and completing the
third yields (MACRO2). Thus every real finite event-free path lifts. H
cannot be revisited after leaving it because G's SCC condensation is
acyclic. Counter reset on reentry cannot therefore mask additional
physical folds.

The lift may still contain impossible paths. It is an exclusion-safe
overapproximation, not a claim that every lifted path is physical.

Complete reconstruction yields 7595 vertices and 7763 edges. Its ONLY
cyclic SCCs are simple cycles of lengths

\begin{verbatim}
24, 104, 28, 28, 28.
\end{verbatim}

The 28-cycles are A phases at counters 0,1,2. The 24- and 104-cycles are
the prior T and standard words. No claim follows from cycle lengths
alone; the words and their nonzero displacements are checked too.

Assign to each condensation component the three weights (one
cyclic-component visit, one noncyclic-vertex update, period):
respectively (1,0,P) on a cyclic component, and (0,1,0) on a noncyclic
singleton. Backward longest-path values are nonnegative and, over every
condensation edge C-\textgreater D, satisfy

\begin{verbatim}
rank(C) >= weight(C)+rank(D)
\end{verbatim}

coordinatewise. Terminal ranks dominate their own weights. Internal
simple-cycle edges keep the same rank. All 7763 lifted edges are
checked. The maxima of the three ranks are

\begin{verbatim}
5, 956, 212.                                          (MACRO3)
\end{verbatim}

Python uses the base's iterative two-pass SCC routine; the independent
native checker reconstructs the lift and uses Tarjan SCCs. Sorted full
edge and rank hashes agree. The source witness fixes the complete edge and rank hashes. Integer inequalities, complete reconstruction and (MACRO2), not a
floating-point solver, carry the proof.

\subsection{Exact physical acceleration on a simple
component}\label{macro-5}

At a cyclic vertex propose its internal periodic word u, of length P and
nonzero displacement d, restoring heading. Let q\_i be its phase read
positions in the actual pose. At q\_i+kd the required old-field
contribution is determined by

\begin{verbatim}
actual required turn bit = 1_B(q_i+kd)
   XOR #{(j,s): q_j-q_i=s d,
                s>0 or (s=0 and j<i), 0<=s<=k} mod 2. (MACRO4)
\end{verbatim}

This computes the literal first mismatch by induction up to that
mismatch; it is not a finite-repeat heuristic. Prior-write changes occur
at finitely many epoch thresholds.

For a finite segment a+jv, 0\textless=j\textless L, its intersection
with a phase ray q+kd is empty, a singleton, or a finite arithmetic
interval in k. Nonparallel directions use an exact determinant. For
parallel directions solve the one-dimensional integer congruence after
dividing by the greatest common divisor. In the closed stride family its
epoch modulus is 1 or 2. This modulus matters: the 24-mode has axis
stride 2 and the 28-mode axis stride 4. Treating every parallel
intersection as a full interval would be wrong.

Split epochs into k=2z and k=2z+1. Each contribution is then an ordinary
finite interval in z; each write threshold is a final interval. Sort
parity-change endpoints separately for the two classes, including z=0
and the final infinite interval. This returns the exact earliest
mismatch, or certifies none. All arithmetic is on integers of polynomial
bit size.

Before a finite mismatch at t=nP+j, append P finite write segments for
the full periods (if n\textgreater0), and j singleton writes for the
partial period. These are XOR contributions; EVERY skipped write is
accounted for. There are at most 2P-1 new terms. Shift the history
vertex around its simple cycle by t mod P. Read the real cell and
execute its actual update. If J10 is true, return the complete
post-state. Otherwise the real edge must leave this simple component:
the unique internal edge has the different predicted label. Continue in
the lifted graph.

If no mismatch exists, the actual infinite continuation equals this
cycle word. A finite field cannot realize the 24- or 28-word forever:
their stabilized first-read classes require respectively two and five
infinite initial black rays, as proved in the prior theorem. Hence only
the standard cycle can give permanent absence. All earlier actual
updates were individually checked or were legal cycle edges, so the
certificate says no J10 at ANY nonnegative time, not only after its
standard boundary.

A component is never revisited after exit. By (MACRO3), one call
therefore uses at most five plans, at most 956 noncyclic updates, at
most five cycle exits, at most 149 warm-up updates, and one final event
update. A safe literal bound is 1111. The sum of visited cycle periods
is at most 212, hence

\begin{verbatim}
new terms <=149+956+2*212+1=1530.
\end{verbatim}

The 2P allocation includes a possible non-event cycle exit. No skipped
write or exit update is omitted. Periodic proposals and inverse-history
computations use a copy until the actual endpoint is determined.

\subsection{A generic finite periodic-plan halo lemma}\label{macro-6}

This lemma concerns ANY supplied periodic turn word of length P
restoring heading with d!=0, from a complete finite source Y=(B,p,h).
Suppose its first mismatch is at time kP+i, position e=p+q\_i+kd, with
phase coordinates relative to p.

Let L be the largest positive translation threshold s with q\_j-q\_i=s
d, or zero if none. Since a period has P unit moves,

\begin{verbatim}
||q_i||_infinity<=P,   L||d||_infinity<=P.
\end{verbatim}

If k\textless=L then
\textbar\textbar e-p\textbar\textbar\_infinity\textless=2P. If
k\textgreater=L+1, all planned previous-write parities at phases i in
epochs k-1 and k have stabilized. If the stable required old-field bit
is 0, mismatch means e is an original black cell. If it is 1, the
preceding occurrence matched, so e-d is an original black cell. Thus in
every case

\begin{verbatim}
e belongs to conv(B union {p})+[-2P,2P]^2.             (MACRO5)
\end{verbatim}

Here B is the FULL source field before the plan, not an unknown active
subset. For k\textgreater=1, any earlier phase point q\_j+ell*d differs
from (ell/k)(q\_i+kd) by at most 2P in infinity norm. For k=0, all reads
are within P of p.~Therefore the entire finite planned read path through
e is inside

\begin{verbatim}
conv(B union {p})+[-4P,4P]^2.                         (MACRO6)
\end{verbatim}

Its ensuing literal update enlarges this by at most one. Unchanged
source cells remain present outside the read footprint; (MACRO6) never
authorizes deleting them.

For radius R=max\_\{b in
B\}\textbar\textbar b-p\textbar\textbar\_infinity (zero for empty B),
(MACRO5)'s more precise alternatives give: either
k\textless=L\textless=P, or
k\textbar\textbar d\textbar\textbar\textless=R+\textbar\textbar d\textbar\textbar+P.
Therefore the first mismatch time is strictly less than

\begin{verbatim}
P(R+P+2).                                           (MACRO7)
\end{verbatim}

These are conditional-on-first-mismatch bounds. A valid infinite plan
has no finite halo bound.

\subsection{All-input J10 geometry, time and bit
complexity}\label{macro-7}

Sum (MACRO6) over all finite periodic portions of one event call. The
period sum is at most 212; all actual literal updates number at most
1111. Convex Minkowski enlargements compose, giving

\begin{verbatim}
H=4*212+1111=1959.
\end{verbatim}

Every intermediate complete state lies in the initial complete convex
hull enlarged by H; in particular its width is at most W0+2H and its
radius about its current head at most that width. Applying (MACRO7) to
every periodic part, using P\textless=104, and including all literal
updates, gives post-event duration at most

\begin{verbatim}
212(W0+2H+106)+1111 =212W0+854199.                   (MACRO8)
\end{verbatim}

After n-1 event calls width is at most W0+2H(n-1). Summing (MACRO8)
gives (MACRO1), since 212H=415308. No theorem bounds n.~In particular an
infinite chain is not classified by the finiteness of each call.

The support of every newly appended segment is made of actual read
positions. Thus all its endpoints lie in the same enlarged initial hull,
even if it later cancels other terms. After n events, all coordinate and
segment-length bit sizes are O(b+log(n+1)); event times also have
O(b+log(n+1)) bits by (MACRO1). The number of terms is at most M+1530n.
Each call performs a fixed number of ray/segment intersections, endpoint
sorts and fixed-length literal or inverse operations. Exact set/list
arithmetic therefore proves poly(M+n+b) bit complexity. Fixed constants
may be large; this is not a benchmark or a claim polynomial in log n.

Together with the already proved H(X) iff finitely many J10 events,
(MACRO1) implies that every hypothetical finite-support HC
counterexample, of EITHER the nonstandard-periodic or nonperiodic type,
has at least n intrinsic J10 events before every time T with (MACRO1)'s
right side \textless=T. In particular its liminf event count divided by
sqrt(T) is at least 1/sqrt(415308). This counts intrinsic events;
interval-contained ten-factor counts differ only in the first nine
boundary positions.

This is a lower rate requirement, not an upper event budget, a
NO-HIGHWAY certificate, or a contradiction. The mass-defect lower bound
and this timing bound can coexist with infinite renewal.
